\section*{Chapter \ref{ch:b1:flowering-plant-organization} --- Functional Organization of a Flowering Plant}

\begin{solution}{exo:b1:flowering-plant-organization:1}
\omterm{def:b1:flowering-plant-organization:organs}{Root}: anchorage, absorption of water and mineral ions. Stem: support
of the \omterm{def:b1:flowering-plant-organization:organs}{leaves} toward the light, conduction between \omterm{def:b1:flowering-plant-organization:organs}{roots} and \omterm{def:b1:flowering-plant-organization:organs}{leaves}.
\omterm{def:b1:flowering-plant-organization:organs}{Leaf}: photosynthesis (and gas exchange, transpiration).
\end{solution}

\begin{solution}{exo:b1:flowering-plant-organization:2}
A module is one node with its \omterm{def:b1:flowering-plant-organization:organs}{leaf} and axillary bud, plus the \omterm{def:b1:flowering-plant-organization:organs}{internode}
below. \omterm{prop:b1:flowering-plant-organization:modular}{Indeterminate growth}: no adult size; the apical and axillary buds
keep adding modules for as long as the plant lives.
\end{solution}

\begin{solution}{exo:b1:flowering-plant-organization:3}
Epidermis: alive; covering, \omterm{def:b1:flowering-plant-organization:tissues}{cuticle}, \omterm{def:b1:flowering-plant-organization:tissues}{stomata}, \omterm{def:b1:flowering-plant-organization:organs}{root hairs}. \omterm{def:b1:flowering-plant-organization:tissues}{Parenchyma}:
alive; photosynthesis, storage. \omterm{def:b1:flowering-plant-organization:tissues}{Collenchyma}: alive; flexible support.
\omterm{def:b1:flowering-plant-organization:tissues}{Sclerenchyma}: dead; rigid support. \omterm{def:b1:flowering-plant-organization:tissues}{Xylem} conducting cells: dead; water
transport. \omterm{def:b1:flowering-plant-organization:tissues}{Phloem} \omterm{def:b1:flowering-plant-organization:tissues}{sieve tubes}: alive (with companion cells); sugar
transport.
\end{solution}

\begin{solution}{exo:b1:flowering-plant-organization:4}
$e^{-0.9} = 0.41$, i.e.\ \qty{41}{\%}; $e^{-3} = 0.05$, i.e.\
\qty{5}{\%}.
\end{solution}

\begin{solution}{exo:b1:flowering-plant-organization:5}
A \omterm{def:b1:flowering-plant-organization:organs}{root} (single central cylinder, \omterm{prop:b1:flowering-plant-organization:stemroot}{endodermis}, \omterm{def:b1:flowering-plant-organization:tissues}{xylem} and \omterm{def:b1:flowering-plant-organization:tissues}{phloem}
alternating on a radius) of a dicot (few \omterm{def:b1:flowering-plant-organization:tissues}{xylem} arms, no central pith).
\end{solution}

\begin{solution}{exo:b1:flowering-plant-organization:6}
A bent beam is stressed most at its surface and not at all along its
axis, so material placed in a peripheral ring resists bending with the
least mass (a tube). A cable under tension is stressed uniformly across
its section and bends freely; a central strand carries the pull while
the cortex stays flexible, and lateral \omterm{def:b1:flowering-plant-organization:organs}{roots} can emerge through it.
\end{solution}

\begin{solution}{exo:b1:flowering-plant-organization:7}
$1 - e^{-2.4} = 0.91$; at $L = 6$, $1 - e^{-3.6} = 0.97$: two more
layers of \omterm{def:b1:flowering-plant-organization:organs}{leaf} add \qty{6}{\%} of the light while costing \qty{50}{\%}
more \omterm{def:b1:flowering-plant-organization:organs}{leaf}.
\end{solution}

\begin{solution}{exo:b1:flowering-plant-organization:8}
$S = \pi d L$: $d = 240/(\pi\times 623\,000) = \qty{1.2e-4}{m}$, about
\qty{0.12}{mm}. $(240 + 400)/5 = 128$.
\end{solution}

\begin{solution}{exo:b1:flowering-plant-organization:9}
\omterm{def:b1:flowering-plant-organization:organs}{Root} surface per mm: $\pi\times 0.4 = \qty{1.26}{mm^2}$. Hairs per mm:
$250\times\pi\times 0.010\times 0.8 = \qty{6.3}{mm^2}$. Total
$7.5/1.26 = 6$: the hairs multiply the surface sixfold.
\end{solution}

\begin{solution}{exo:b1:flowering-plant-organization:10}
In the forest the lower axillary buds are shaded; the modules they
would make cannot pay for themselves, so they stay dormant or their
branches die, and growth is allocated to the apical bud racing upward
toward the light. In the field every bud is lit and every branch pays,
so the crown fills out low. The cost: a tall thin trunk is a large
investment in unproductive wood, vulnerable to wind; the gain: the plant
matches its form to where the light is, without any fixed plan.
\end{solution}

\begin{solution}{exo:b1:flowering-plant-organization:11}
The pit and hairs hold a layer of still, humid air above the pore; the
vapour gradient from the air spaces to the moving air is spread over a
longer path and transpiration falls. Carbon dioxide must diffuse in
along the same longer path, so uptake falls too, though less, since the
$\mathrm{CO_2}$ gradient is set by the atmosphere and the \omterm{def:b1:flowering-plant-organization:organs}{leaf}'s
consumption. In a desert water is the limiting resource and the trade
saves the plant's life; in a rain forest water is free and carbon gain
is what competition rewards.
\end{solution}

\begin{solution}{exo:b1:flowering-plant-organization:12}
The plant has no regulated \omterm{def:b1:mammal-organization:compartments}{extracellular fluid} held at constant
temperature and composition between its cells and the outside: its
cells meet soil water, air and sunlight directly, and its temperature
is the air's. Each cell therefore regulates its own contents (ions,
water, pH) across its membrane and wall; the plant regulates its
exchanges at the surface (\omterm{def:b1:flowering-plant-organization:tissues}{stomata} open and close, the \omterm{prop:b1:flowering-plant-organization:stemroot}{endodermis}
selects what enters the \omterm{def:b1:flowering-plant-organization:tissues}{xylem}) and adjusts its body by growth (more
\omterm{def:b1:flowering-plant-organization:organs}{roots} in dry soil, more \omterm{def:b1:flowering-plant-organization:organs}{leaves} in shade). It regulates, but the
regulation is at the surfaces and in the form, not in a private fluid.
\end{solution}

\begin{solution}{pb:b1:flowering-plant-organization:1}
\textbf{1.} $5\times 100 = \qty{500}{m^2}$.
\textbf{2.} $e^{-2.5} = 0.082$ reaches the ground; \qty{92}{\%}
intercepted.
\textbf{3.} $L = 2$: $1 - e^{-1} = \qty{63}{\%}$; $L = 8$: $1 -
e^{-4} = \qty{98}{\%}$.
\textbf{4.} From 2 to 5 the crown gains \qty{29}{\%} of the light for
\qty{300}{m^2} of \omterm{def:b1:flowering-plant-organization:organs}{leaf}; from 5 to 8 it would gain \qty{6}{\%} for
another \qty{300}{m^2} that cost as much to build and to keep supplied
with water: the extra \omterm{def:b1:flowering-plant-organization:organs}{leaves} do not pay.
\textbf{5.} $500/0.005 = \num{100000}$ \omterm{def:b1:flowering-plant-organization:organs}{leaves}.
\textbf{6.} $500\times 10^6\,\mathrm{mm^2}\times 150 = \num{7.5e10}$
\omterm{def:b1:flowering-plant-organization:tissues}{stomata}.
\textbf{7.} $\qty{50}{\micro m^2}$; per mm$^2$: $150\times 50\times
10^{-6} = \qty{7.5e-3}{mm^2}$, i.e.\ \qty{0.75}{\%} of the surface.
\textbf{8.} $500\times 4\times 10^{-6} = \qty{2e-3}{mol/s}$, i.e.\
\qty{88}{mg/s}.
\textbf{9.} $0.088\times 43\,200 = \qty{3.8}{kg}$ of $\mathrm{CO_2}$,
containing $3.8\times 12/44 = \qty{1.04}{kg}$ of carbon.
\textbf{10.} $150\times 1.04 = \qty{156}{kg}$ of carbon; half,
\qty{78}{kg}, into wood at \qty{50}{\%} carbon: \qty{156}{kg} of wood.
\textbf{11.} $500\times 10^{-3} = \qty{0.5}{mol/s}$, i.e.\ \qty{9}{g/s}.
\textbf{12.} $9\times 43\,200 = \qty{389}{kg}$, about \qty{390}{L}.
\textbf{13.} $0.5/(2\times 10^{-3}) = 250$ water molecules per
$\mathrm{CO_2}$. The pore that lets $\mathrm{CO_2}$ in lets water out;
the inside air is saturated while $\mathrm{CO_2}$ is only
\qty{0.04}{\%} of the outside air, so the outward gradient of vapour is
far steeper than the inward gradient of $\mathrm{CO_2}$.
\textbf{14.} $\qty{2}{mm}\times\qty{100}{m^2} = \qty{200}{L}$: half a
day.
\textbf{15.} \qty{100}{m^3}; $5\times 100 = \qty{500}{km}$.
\textbf{16.} $\pi\times 0.5\times 10^{-3}\times 5\times 10^5 =
\qty{785}{m^2}$.
\textbf{17.} $200\times 5\times 10^8\,\mathrm{mm} = \num{1e11}$ hairs.
\textbf{18.} Each $\pi\times 10^{-5}\times 5\times 10^{-4} =
\qty{1.57e-8}{m^2}$; total \qty{1570}{m^2}.
\textbf{19.} $785 + 1570 = \qty{2355}{m^2}$, $4.7$ times the \omterm{def:b1:flowering-plant-organization:organs}{leaf} area.
\textbf{20.} Volume within \qty{1}{mm}: $\pi\times(10^{-3})^2\times 5
\times 10^5 = \qty{1.6}{m^3}$, i.e.\ \qty{1.6}{\%} of the soil.
\textbf{21.} Rye: $\pi\times 10^{-6}\times 6.2\times 10^5 =
\qty{1.95}{m^3}$ of the \qty{0.05}{m^3} box, i.e.\ \qty{3900}{\%}: every
point of the soil is within a millimetre of forty \omterm{def:b1:flowering-plant-organization:organs}{roots}. The rye
explores its soil more than two thousand times more thoroughly ---
a crop plant on a four-month budget against a tree exploring a hundred
cubic metres.
\textbf{22.} $1000 + 2355 = \qty{3355}{m^2}$; $\qty{34}{m^2}$ per
square metre of ground.
\textbf{23.} Human: $130/2 = 65$; the oak's ratio is half that of the
human's folded surfaces, without any folding.
\textbf{24.} The oak adds modules, \omterm{def:b1:flowering-plant-organization:organs}{leaves} and \omterm{def:b1:flowering-plant-organization:organs}{roots} every year, so its
surfaces grow with its age; the mammal's surfaces are fixed \omterm{def:b1:mammal-organization:organ}{organs},
complete at adult size, and the only way to raise its ratio was to fold
them during development.
\textbf{25.} About \qty{34}{m^2} of \omterm{prop:b1:organism-environment:surfaces}{exchange surface} per square metre
of ground: \qty{10}{m^2} of \omterm{def:b1:flowering-plant-organization:organs}{leaf} (both faces) above and \qty{24}{m^2}
of \omterm{def:b1:flowering-plant-organization:organs}{root} and \omterm{def:b1:flowering-plant-organization:organs}{root hair} below.
\end{solution}
